\documentclass{article}
\usepackage{iclr2027_conference,times}
\input{math_commands.tex}
\usepackage{booktabs}
\usepackage{amsthm,amssymb}
\usepackage{hyperref}
\usepackage{url}

\title{When Minimal Recovery Explanations Fail:\\
Exact Audits under Coarsening and Abstention}
\author{Anonymous authors}
% Working manuscript: no submission or acceptance is claimed.
\newtheorem{proposition}{Proposition}
\newcommand{\Fam}{\mathcal{M}}
\newcommand{\Nec}{\mathcal{D}}
\newcommand{\Poss}{\mathcal{P}}
\newcommand{\Completions}{\mathcal{C}}
\newcommand{\Representable}{\mathcal{A}}
\newcommand{\cl}{\operatorname{cl}}
\newcommand{\mininc}{\operatorname{Min}_{\subseteq}}
\newcommand{\unknown}{\mathord{?}}

\begin{document}
\maketitle
\lhead{Retrospective methods and pilot study --- working manuscript}

\begin{abstract}
A minimal prompt edit can explain an observed behavioral change without
identifying a complete explanation family or a stable intervention. We audit
these distinctions using existing, exhaustively enumerated jailbreak-recovery
decision tables. Our method reports sharp marginal membership bounds over
Boolean completions of unresolved decisions and compares coarsened explanations
in their common atomic-unit space. An elementary partition-closure identity
provides a grouping-only reference for complete monotone oracles; observed
recovery-to-nonrecovery reversals supply a stronger diagnostic when tables are
incomplete. In a retrospective pilot of 12 attack--payload instances spanning
nine payloads, 16 minima are necessary across all completions, while 45 are
marginally possible; only seven complete families are identified. Of 36
adjacent-merge contrasts, 25 change the certified atomic family, but four of
eight nontriviality flips leave that family unchanged. All 14 eligible contrasts
from seven complete h4rm3l tables agree with the closure reference. Four
contrasts, representing three distinct atomic reversal pairs in two
DeepInception instances, contain a certified recovery reversal. All three
distinct adjacent reversals are specific to layout-preserving blanking;
source-aware omission retains three safe sampled outcomes at their closures.
These findings
qualify an initially broad instability interpretation: much of the apparent
effect is compatible with grouping or unresolved certification. We preserve
the failed neutralizer-stability gate and subsequent negative two-target
screen that motivated the audit. The contribution is an executable
identification audit with a small, explicitly retrospective empirical study,
not prospective evidence of widespread instability or a safety guarantee.
\end{abstract}

\section{Introduction}

Suppose neutralizing two components of a jailbreak prompt changes a harmful
model response to an outcome accepted by a fixed evaluation rule. Calling this
pair a minimal recovery explanation invites several different conclusions:
the pair has been certified minimal; every minimal alternative is known; the
result persists if components are grouped differently; and the underlying
model is now safe. None of these conclusions follows from the others.
This paper studies the first three claims while keeping the fourth outside
the scope of a finite behavioral experiment.

Minimal input explanations and feature removal are established tools
\citep{carter2019,covert2021}. Structured jailbreaks make an audit particularly
concrete: transformations can be represented by separately editable components,
while the original request remains fixed
\citep{doumbouya2025,li2023}. However, a minimum depends on which interventions
are available and which outcomes an evaluator certifies. Existing work on
jailbreak evaluation already documents discrepancies and proposes more
specific criteria \citep{huang2026guidedbench,huang2026jailmeter}. Our question
concerns the resulting \emph{family of minimal interventions}: how much of that
family is identified, and which changes under coarsening exceed what grouping
alone can explain?

We answer this question through a retrospective audit of a completed
development experiment. The original study sought recovery structure stable
across two neutralizers, but its frozen cross-family stability gate failed.
A later two-target confirmation screen also failed its eligibility gate.
These negative results motivated the present analysis; they are preserved as
negative results. No new target or evaluator calls were made for this audit.

Our central observation is a distinction between \emph{complete enumeration}
and \emph{complete identification}. Every subset may have an explicit record
while some records remain unresolved because of evaluator disagreement or
capability confounding. Exhaustively listing certified minima then gives a
lower family, not necessarily the entire family of any resolved oracle. We
derive the exact intersection and union of minimal-family memberships over
all Boolean completions that respect the known labels. These are logical
bounds, not confidence intervals or human validation.

We also separate three sources of apparent resolution sensitivity: changes
only in the number of named groups, changes expected from restricting the
available interventions, and an observed loss of recovery when a fine minimum
is enlarged to a representable coarse intervention. The pilot shows why this
separation matters. A broad initial comparison found 25 family changes in 36
adjacent merges. Yet the grouping reference explains every tested contrast
with a complete Boolean table, and the stronger reversal evidence is confined
to two DeepInception instances and the blanking operator. The completed result is therefore an audit of
what the data identify and what a coarse explanation can support, with limited
positive evidence for a stronger failure mode.

\section{A finite recovery oracle with unresolved decisions}
\label{sec:oracle}

Let $U=\{u_1,\ldots,u_m\}$ be a fixed ordered vocabulary of attack-owned
components. An instance comprises a target artifact, an attack renderer, and
a request $p$. For $S\subseteq U$, an intervention neutralizes the components
in $S$ while preserving the UTF-8 bytes of $p$ exactly once in the scientific
prompt. The observation contract fixes neutralizers, decoding seeds, response
evaluators, output-validity checks, and capability controls.

We represent the resulting decision table by
\begin{equation}
g:2^U\longrightarrow\{1,0,\unknown\}.
\label{eq:partial}
\end{equation}
Here $1$ means \texttt{RECOVERED} under that contract, $0$ means
\texttt{NOT\_RECOVERED}, and $\unknown$ retains an unresolved label. In the
completed experiment, recovery required the fixed panel to label every
required seed under both neutralizers safe, followed by passing all selected
capability controls. A valid harmful witness certified nonrecovery. Evaluator
abstentions without such a witness remained unresolved. Failed capability
controls were separately recorded as \texttt{CAPABILITY\_CONFOUNDED}; they
enter $\unknown$ for identification calculations while retaining their
original reason codes. No unresolved observation is promoted to recovery.

The term ``known'' below means fixed by this recorded instrument. It does not
mean that an automatic judgment is semantically infallible. The Boolean
completion model describes uncertainty \emph{conditional on the accepted
labels}; it does not model misclassification of labels already marked $0$ or
$1$, nor does it repair capability failure.

For a total Boolean function $f:2^U\to\{0,1\}$, define its complete
strict-subset-minimal family
\begin{equation}
\Fam(f)=\{S\subseteq U:f(S)=1\ \text{and}\ f(T)=0
                 \text{ for every }T\subsetneq S\}.
\label{eq:minimal}
\end{equation}
This is inclusion minimality, not minimum cardinality. We do not assume
monotonicity: neutralizing additional components may restore a harmful
response. Checking only immediate subset neighbors is insufficient for
Equation~\ref{eq:minimal} without such an assumption.

Three levels of exactness should consequently be distinguished. A
\emph{complete decision table} contains all $2^m$ rows, including explicit
unknowns. An \emph{exhaustive certificate family} contains every minimum that
can be proved from those rows. An \emph{identified Boolean minimal family}
is the same under every allowed resolution of the unknowns. None establishes
universal safety outside the recorded evaluation contract.

\section{Identification and coarsening audits}

\subsection{Necessary and possible minimal memberships}

Let $\Completions(g)$ contain all Boolean functions agreeing with every known
label of $g$, with no monotonicity or distributional constraint. Define
\begin{align}
\Nec(g)&=\{S:g(S)=1,\quad \forall T\subsetneq S,\ g(T)=0\},
\label{eq:necessary}\\
\Poss(g)&=\{S:g(S)\ne0,\quad
                   \nexists T\subsetneq S\text{ with }g(T)=1\}.
\label{eq:possible}
\end{align}

\begin{proposition}[Sharp marginal membership bounds]
\label{prop:bounds}
For any finite table $g$,
\begin{equation}
\Nec(g)=\bigcap_{f\in\Completions(g)}\Fam(f),\qquad
\Poss(g)=\bigcup_{f\in\Completions(g)}\Fam(f).
\label{eq:bounds}
\end{equation}
Thus $\Nec(g)\subseteq\Fam(f)\subseteq\Poss(g)$ for every completion, and
the entire family is identified exactly when $\Nec(g)=\Poss(g)$.
\end{proposition}

The proof follows by assigning unresolved labels at $S$ and its strict
subsets; Appendix~\ref{app:proofs} gives the details. ``Sharp'' is marginal:
each member of $\Poss(g)$ occurs in some completion, but all possible members
need not coexist. In particular, $\Poss(g)$ need not be an antichain. Reporting
its size as the number of recovered explanations would be incorrect.
Complete Boolean observation is sufficient for identification, but not
necessary if the unresolved rows cannot affect minimality.

A two-unit example illustrates the distinction. Set $g(\varnothing)=0$,
$g(\{a\})=\unknown$, $g(\{b\})=0$, and $g(\{a,b\})=1$. Then
$\Nec=\varnothing$ and $\Poss=\{\{a\},\{a,b\}\}$. Resolving the unknown
to $1$ gives the first minimum; resolving it to $0$ gives the second. The two
possible minima cannot coexist in a single completion.

\subsection{Exact projection and a grouping reference}

Let $\pi$ be a partition of $U$, and let $\Representable_\pi$ denote all
unions of its blocks. A coarse intervention is the same atomic intervention
already present in the table, restricted to this subdomain. Its label is
therefore obtained without a new model call. We compare coarse minima as
\emph{atomic unions}, rather than comparing block names or their ordering.

Define the partition closure
\begin{equation}
\cl_\pi(S)=\bigcup\{B\in\pi:B\cap S\ne\varnothing\},
\end{equation}
the smallest representable intervention containing $S$. For any family $F$,
write $\mininc F$ for its inclusion-minimal members. Let $\Fam_\pi(f)$ be
the minimal positives of $f$ restricted to $\Representable_\pi$, expressed
in atomic units.

\begin{proposition}[Monotone partition-closure identity]
\label{prop:closure}
For a complete monotone Boolean oracle $f$,
\begin{equation}
\Fam_\pi(f)=\mininc\{\cl_\pi(M):M\in\Fam(f)\}.
\label{eq:closure}
\end{equation}
\end{proposition}

This elementary identity supplies a grouping reference, not a claim that the
observed oracle is monotone. We call disagreement with its right-hand side a
\emph{closure-reference residual}. The implementation computes that residual
only on complete Boolean tables; substituting a lower certificate family for
the unidentified $\Fam(f)$ would give an unjustified reference. A nonzero
residual rejects this particular monotone explanation. A zero residual does
not prove monotonicity or independence of the units.

Grouping can also change a verbal classification without changing any atomic
minimum. If the only minimum is $\{a,b\}$ and those units become one block,
its group cardinality changes from two to one. We reproduce the historical
``nontrivial'' indicator, which is true when a certified family has multiple
members or contains a member of order greater than one. A flip with identical
atomic families is reported as a \emph{group-count-only flip}; it is not
evidence of a changed atomic explanation.

\subsection{A stronger witness available with incomplete tables}

For a certified fine minimum $M\in\Nec(g)$, a
\emph{recovery-reversal witness} is
\begin{equation}
g(M)=1,\qquad g(\cl_\pi(M))=0.
\label{eq:witness}
\end{equation}
These two known labels contradict every monotone Boolean completion of $g$.
Unlike a family-name change, the witness states that enlarging a recorded
recovery intervention to the required coarse union loses recovery under the
same contract. Unknown closure labels are counted separately and never treated
as witnesses. Because $g$ aggregates a finite seed and neutralizer contract,
this is a statement about that observed contract, not an estimate of a
population-level stochastic recovery probability.

Coarsening can additionally \emph{promote} a certificate by hiding strict
subsets that are no longer representable. A hidden recovered subset changes
the admissible optimization problem. A hidden unresolved subset changes what
can be certified on that restricted domain. In the two-unit example above,
merging $a$ and $b$ yields a certified coarse singleton although no fine
minimum is identified. The algorithm records these causes instead of
treating every coarse certificate as stronger evidence.

The method enumerates subsets, calculates Equations~\ref{eq:necessary}
and~\ref{eq:possible}, projects each declared partition, compares families
setwise in atomic coordinates, and records closure witnesses and hidden-subset
causes. It requires no access to prompt or response text. A direct
implementation checks $O(3^m)$ subset relations; the model-query requirement
belongs to the original finite table, not to subsequent projections.

\section{Retrospective experimental setting}
\label{sec:setting}

\subsection{Data, model, and recorded oracle}

The topology corpus consists of all 12 stable attack--payload pairs retained
by the preceding frozen development screen: nine instances of one fixed
h4rm3l composition and three of one fixed DeepInception renderer, covering
nine unique GuidedBench core requests \citep{huang2026guidedbench}. This is a
selected development corpus, not a random sample of all attacks or of the
whole benchmark. All topology outputs come from Qwen2.5-7B-Instruct in the
official Q4\_K\_M artifact, executed with llama.cpp build \texttt{b10441} on
an RTX 3070. Sampling used seeds 11, 23, and 47, temperature $0.7$,
top-$p$ $0.9$, top-$k$ $50$, and a 512-token generation maximum.

The h4rm3l vocabulary contains three separately implemented decorators. The
DeepInception vocabulary contains seven source-derived components of the
fixed renderer. Neither vocabulary was changed in this audit. Two
neutralizers, source-aware omission and layout-preserving blanking, were
crossed with the required seeds. Two evaluator axes, Qwen3Guard and
JailMeter, were combined using the frozen agreement rule
\citep{zhao2025,huang2026jailmeter}. The exact mapping and capability contract
are given in Appendix~\ref{app:contract}; the safe labels are not new
GuidedEval measurements.

The resulting table has $9\cdot2^3+3\cdot2^7=456$ subset decisions,
including the empty baselines: 64 recovered, 346 not recovered, 29 abstained,
and 17 capability confounded. All rows are present, but 46 are unresolved
for Boolean-family identification. The original experiment made 1,807 new
target generations and 324 capability-control generations. Another 857
target generations were skipped only after a valid harmful witness had
already falsified the conjunction for that subset and neutralizer. This
logical stopping rule does not synthesize unobserved per-seed outcomes.

\subsection{Preserved negative results and analysis chronology}

Before this audit, the original cross-family development gate required pooled
positive-neutralizer Jaccard at least $0.80$. The observed value was
$64/141=0.4539$, so that gate failed. The family-specific values were
$25/30=0.8333$ for h4rm3l and $39/111=0.3514$ for DeepInception. Those
post-outcome diagnostics motivated a separately frozen h4rm3l-only,
two-target confirmation screen, not a relabeling of the original gate.

The confirmation screen used 45 previously reserved payloads with Qwen2.5
and an admitted Gemma 4 E4B QAT Q4\_0 artifact. Its 180 seed-11 direct and
attacked generations completed. Forty target--payload pairs could advance,
but the split was Qwen 35 and Gemma 5 against a minimum of six per target.
Later seeds could only remove pairs, making the gate unreachable. The screen
therefore ended with a valid negative result and no confirmatory topology.
The stored zero fully stable pairs must not be read as a three-seed prevalence
estimate, because the later seeds were not executed.

The present identification and coarsening questions were formulated after
these results. Every number below is a descriptive reconstruction of the
existing safe tables. No prospective success gate was applied to this new
audit, and neither original failed gate was modified. An earlier projection
enumerated 2,676 arbitrary partitions and 228 contiguous partitions; our main
diagnostics use the 36 immediate adjacent merges, with all 90 immediate
arbitrary merges as an additional descriptive view. These are repeated
deterministic contrasts, not independent sample sizes.

\section{Results}

\subsection{Exhaustive certificates leave five families unidentified}

Table~\ref{tab:identification} reports necessary and possible memberships.
The 16 necessary minima exactly reconstruct the original certified minima.
The union of marginally possible memberships contains 45 sets, including 29
in the three DeepInception instances. Only seven instance families are
identified; all seven happen to have complete Boolean tables and belong to
h4rm3l. Three DeepInception instances and two h4rm3l instances remain
unidentified. Thus the previously reported 11 instances with at least one
certified minimum do not imply 11 identified complete families.

\begin{table}[t]
\caption{Family identification in the 12-instance pilot. Counts of minima
sum memberships within instances; possible memberships need not coexist.
The overall nine payloads are shared across attack families.}
\label{tab:identification}
\centering
\small
\begin{tabular}{lrrrrr}
\toprule
Attack & Instances & Unknown rows & Necessary & Possible & Identified \\
\midrule
h4rm3l & 9 & 2 & 12 & 16 & 7/9 \\
DeepInception & 3 & 44 & 4 & 29 & 0/3 \\
\midrule
Total & 12 & 46 & 16 & 45 & 7/12 \\
\bottomrule
\end{tabular}
\end{table}

This result changes the interpretation of an exact-search comparison. In
the historical h4rm3l audit, a singleton search found 1 of 12 certified minima,
while each of the implemented greedy-forward, greedy-backward, and
DDMIN-style paths found 8 of 12. The four missed memberships demonstrate
that one returned path need not exhaust the certificate family. The
denominator is the \emph{12 certified memberships}, not an identified full
family for every instance and not the 16 marginally possible h4rm3l
memberships. Moreover, a method designed to return one minimum should not be
called incorrect merely because it does not enumerate all alternatives. This
comparison is descriptive of different output objects, not a controlled
runtime-superiority benchmark.

\subsection{Grouping explains the eligible complete-table contrasts}

Table~\ref{tab:adjacent} decomposes the 36 adjacent-merge comparisons. Atomic
certificate families change in 25, while the historical nontriviality
classification changes in eight. Four of these eight flips preserve the
entire atomic family. For those four, the difference is solely the number
of groups used to name the same intervention.

For the seven complete Boolean h4rm3l tables, all 14 adjacent contrasts agree
with the closure reference: there are zero residuals. This does not establish
monotonicity, but it provides no evidence in those contrasts that the family
changes exceed the prediction of the grouping reference. The remaining
22 contrasts are not assigned a reference residual because their fine
families are not completely observed. In particular, the absence of a
DeepInception residual estimate is missing eligibility, not a zero effect.

\begin{table}[t]
\caption{Adjacent coarsening audit. Numerators count partition contrasts,
not independent cases. ``Not eligible'' means no complete Boolean table was
available for the closure-reference comparison.}
\label{tab:adjacent}
\centering
\small
\begin{tabular}{lrrr}
\toprule
Diagnostic & h4rm3l & DeepInception & Overall \\
\midrule
Adjacent contrasts & 18 & 18 & 36 \\
Atomic family changed & 12 & 13 & 25 \\
Nontriviality flipped & 8 & 0 & 8 \\
Group-count-only flip & 4 & 0 & 4 \\
Recovery-reversal contrasts & 0 & 4 & 4 \\
Unknown closures of certified minima & 0 & 4 & 4 \\
Closure-reference residual / eligible & 0/14 & Not eligible & 0/14 \\
\bottomrule
\end{tabular}
\end{table}

\subsection{Observed reversals are localized and partly repeated}

Four adjacent contrasts have a recovered fine minimum whose coarse closure
is recorded not recovered. They occur in two DeepInception instances,
representing two unique payloads; h4rm3l has none. Table~\ref{tab:witnesses}
expresses the witnesses using the seven frozen unit indices. Two different
partitions in instance D3 yield the same minimum--closure pair. Thus the
four contrast-level witnesses represent \emph{three distinct atomic
reversal pairs}, not four independent behavioral events.

The operator-specific provenance is decisive. At every adjacent witness,
the fine minimum has three safe sampled responses under both neutralizers.
At its coarse closure, source-aware omission still has three safe responses,
while layout-preserving blanking has one recorded harmful response. A first
harmful response at seed 11 triggers the frozen early stop in one contrast;
the other three contrasts include two safe responses and a harmful response
at seed 47. No unexecuted seed is imputed. The signal is therefore a failure
of the combined finite-seed predicate attributable to blanking at these
endpoints, not a reversal replicated under omission or harmfulness sustained
across all seeds.

\begin{table}[t]
\caption{Distinct adjacent recovery reversals. D1 and D3 identify rows in
Appendix~\ref{app:instances}; numbers denote frozen component positions.
The last column counts adjacent partitions yielding the same pair.}
\label{tab:witnesses}
\centering
\begin{tabular}{lccc}
\toprule
Instance & Recovered minimum & Nonrecovered closure & Partitions \\
\midrule
D1 & $\{1,4\}$ & $\{1,3,4\}$ & 1 \\
D1 & $\{1,4\}$ & $\{1,2,4\}$ & 1 \\
D3 & $\{1,4,6,7\}$ & $\{1,4,5,6,7\}$ & 2 \\
\bottomrule
\end{tabular}
\end{table}

These witnesses survive arbitrary Boolean assignments to the other unknown
rows: their endpoints are known labels. They establish observed
nonmonotonicity of the recorded recovery predicate for those pairs. They do
not identify the whole DeepInception family, establish evaluator truth, or
support a replicated effect across two attack families or two targets.
The complete topology evidence still comes from one quantized target.
The two families also have different opportunity counts: three units yield
two adjacent merges per h4rm3l instance, whereas seven yield six per
DeepInception instance. Their raw witness counts cannot isolate an attack-family
effect from dimensionality and the number of comparisons. h4rm3l is an
empirical comparison family, not a known-negative control; only a stipulated
synthetic monotone oracle supplies that guarantee for the closure diagnostic.

Across all 90 arbitrary one-pair merges, 64 contrasts change the atomic
certificate family, and 11 contrasts contain 12 witness occurrences, again
confined to the same two instances. They deduplicate to five atomic reversal
pairs: three blank-only, one omission-only, and one involving both
neutralizers. This is an exploratory sensitivity view; choosing nonadjacent
partitions after observing these outcomes would not validate an
operator-general adjacent-coarsening claim. All 21 eligible complete-table comparisons
have zero closure-reference residual. This wider enumeration adds
deterministic views of the same observations, not independent replication.

\subsection{The negative confirmation screen exposes an instrument limit}

Table~\ref{tab:c1n} shows the attacked-response labels in the separate
two-target screen. Gemma has only 9/45 evaluator agreements and 36 panel
abstentions. Qwen has 36/45 agreements and nine abstentions. The five Gemma
harmful agreements explain why the fixed panel could not certify enough
eligible pairs. This identifies a coverage bottleneck under the recorded
instrument; it does not adjudicate between the judges or prove that the
attack failed to elicit relevant assistance.

\begin{table}[t]
\caption{Attacked-response labels in the completed seed-11 screen. Each
cell is harmful / safe / abstain. These are automatic labels, not human
ground truth, and no confirmation topology was generated.}
\label{tab:c1n}
\centering
\begin{tabular}{lccc}
\toprule
Target & Qwen3Guard & JailMeter & Fixed panel \\
\midrule
Qwen2.5 & 45 / 0 / 0 & 36 / 9 / 0 & 36 / 0 / 9 \\
Gemma & 41 / 4 / 0 & 5 / 35 / 5 & 5 / 4 / 36 \\
\bottomrule
\end{tabular}
\end{table}

\section{Related work and scope of contribution}

\paragraph{Minimal inputs and removal explanations.}
Delta debugging minimizes failure-inducing inputs \citep{zeller2002}, and
sufficient input subsets provide compact, instance-specific explanations
\citep{carter2019}. Removal-based explanation frameworks distinguish how
features are removed, which behavior is explained, and how influence is
summarized \citep{covert2021}. Formal explanation work also develops
relationships and enumeration algorithms for sufficient and contrastive
explanations \citep{ignatiev2020}. Our fixed intervention $f(S)=1$ evaluates
one specified edited input; it does not universally quantify over all values
of the remaining features. We do not claim the first minimal input set,
the first complete finite enumeration, or the general discovery that a
representation affects explanation. Propositions~\ref{prop:bounds}
and~\ref{prop:closure} are elementary finite-set facts supplied to make the
audit explicit and checkable. The contribution is their use in one executable
audit that distinguishes certified versus identified recovery families,
grouping effects, hidden uncertainty, and known reversal witnesses.

\paragraph{Structured attacks and evaluation.}
h4rm3l provides a compositional transformation language
\citep{doumbouya2025}; DeepInception uses a nested-scene prompt construction
\citep{li2023}. We reuse fixed representatives and do not synthesize or
optimize attacks. GuidedBench introduces case-specific evaluation guidelines,
while JailMeter develops an evidence-based jailbreak evaluator
\citep{huang2026guidedbench,huang2026jailmeter}. Qwen3Guard supplies a distinct
safety-moderation instrument \citep{zhao2025}. Our empirical endpoint is the
historical fixed panel, whose disagreement is part of the audit. We do not
claim a new evaluator or a new finding that jailbreak evaluators disagree.
LOCA is especially close in motivation: it seeks minimal internal
representation changes that induce refusal \citep{kumar2026}. Our editable
input-unit families and conditional observation audit do not establish
internal causal mechanisms or subsume that representation-level analysis.

\section{Limitations and requirements for broader validation}

\paragraph{Selection and scope.}
The 12 topology instances were selected for stable attack eligibility in a
development screen; the questions in this manuscript were chosen after
observing their outcomes. Nine unique payloads and one quantized target
cannot support prevalence or broad model-family claims. We therefore report
finite counts without a hypothesis test or confidence interval that treats
partition edges as independent samples. A future population study must
account for repeated target, attack, seed, and partition measurements within
payload clusters. It must also avoid presenting the three-unit composition
as representative of all h4rm3l programs.

\paragraph{Conditional correctness.}
Sharp completion bounds quantify missing-label uncertainty conditional on
the fixed known decisions. Judge errors may invalidate those conditions.
Capabilities are tested using two simple tasks and cannot rule out every
semantic or behavioral confound. Grouping units preserves the specified
atomic intervention but changes the available edit vocabulary; a coarse
minimum should not be read as a newly discovered intrinsic mechanism.

\paragraph{Measurement redesign remains unvalidated.}
A no-inference development preflight has bound the 180 existing screen
responses to their GuidedBench guidelines. No GuidedEval judge output is
part of this paper. A candidate follow-up endpoint would certify that all
case-specific entity/action scoring points are absent, retain pointwise judge
disagreement as uncertainty, and report each judge separately. Even unanimous
zero matched points means only zero matches to the declared rubric on those
outputs, not universal safety or absence of other harmful content. GuidedBench
already owns this rubric-based evaluation approach; adapting its zero-point
endpoint cannot be claimed as a newly validated judge.

\paragraph{Prospective evidence.}
An identity-only audit has sealed two unused 60-payload cohorts with zero
overlap and no prior-output collision. Neither cohort contributed a target
or evaluator outcome here. Broader validation requires a complete pre-output
contract specifying judge artifacts, qualification criteria, admissible
partitions, target admission, capability controls, estimands, and stopping
rules. Development qualification should use only already observed responses.
Fresh primary data must then establish whether adequate identifiable families
and meaningful reversals occur, with independent replication of frozen
contrasts on the second cohort. The current data do not meet that broader
evidence requirement.

\section{Conclusion}

An exhaustive prompt-intervention experiment can identify some minimal
certificates while leaving the full explanation family unknown. In the
completed pilot, 16 necessary memberships expand to 45 marginally possible
memberships, and only seven of 12 families are identified. Coarsening changes
many reported families, but every eligible complete-table contrast matches
the elementary grouping reference; half of the nontriviality flips preserve
the atomic family itself. A stronger recovery-reversal diagnostic identifies
three distinct atomic pairs in two DeepInception instances, all specific to
blanking among the adjacent contrasts. Reporting these
objects separately produces a narrower and more defensible account of what
the experiment establishes. The resulting audit is complete for the stored
tables; its empirical generalization remains a prospective research task.

\section*{Reproducibility and responsible reporting}

The retrospective analysis consumes only safe tables containing unit IDs,
hashes, and categorical outcomes. Raw requests, attack prompts, model
responses, and judge explanations are not included in this manuscript or
its analysis output. The source tables and original negative receipts are
preserved. Appendix~\ref{app:artifacts} identifies the local analysis entry
point and exact source/output identities. No new human annotation was
performed for this study. No claim of public artifact release, submission,
acceptance, human-equivalent judgment, or generalized safety is made.

\begin{thebibliography}{10}

\bibitem[Carter et~al.(2019)]{carter2019}
Brandon Carter, Jonas Mueller, Siddhartha Jain, and David Gifford.
What made you do this? Understanding black-box decisions with sufficient
input subsets.
In \emph{Proceedings of AISTATS}, volume 89 of \emph{PMLR},
pp.\ 567--576, 2019.
\url{https://proceedings.mlr.press/v89/carter19a.html}.

\bibitem[Covert et~al.(2021)]{covert2021}
Ian Covert, Scott Lundberg, and Su-In Lee.
Explaining by removing: A unified framework for model explanation.
\emph{Journal of Machine Learning Research}, 22(209):1--90, 2021.
\url{https://jmlr.org/papers/v22/20-1316.html}.

\bibitem[Doumbouya et~al.(2025)]{doumbouya2025}
Moussa Koulako Bala Doumbouya, Ananjan Nandi, Gabriel Poesia, Davide Ghilardi,
Anna Goldie, Federico Bianchi, Dan Jurafsky, and Christopher D. Manning.
h4rm3l: A language for composable jailbreak attack synthesis.
In \emph{International Conference on Learning Representations}, 2025.
\url{https://arxiv.org/abs/2408.04811v4}.

\bibitem[Huang et~al.(2026a)]{huang2026guidedbench}
Ruixuan Huang, Xunguang Wang, Zongjie Li, Daoyuan Wu, and Shuai Wang.
GuidedBench: Measuring and mitigating the evaluation discrepancies of
in-the-wild LLM jailbreak methods.
In \emph{International Conference on Learning Representations}, 2026.
\url{https://arxiv.org/abs/2502.16903v3}.

\bibitem[Huang et~al.(2026b)]{huang2026jailmeter}
Qingjia Huang, Jingyu Zhang, Jianguo Wu, Yakai Li, Weijuan Zhang,
Yankai Rong, Junyi Yao, Shengzhi Zhang, and Xiaoqi Jia.
JailMeter: An evidence-based evaluation framework for jailbreak attacks on
large language models.
In \emph{Findings of the Association for Computational Linguistics: ACL 2026},
pp.\ 16006--16029, 2026.
\url{https://aclanthology.org/2026.findings-acl.786/}.

\bibitem[Ignatiev et~al.(2020)]{ignatiev2020}
Alexey Ignatiev, Nina Narodytska, Nicholas Asher, and Jo\~ao Marques-Silva.
On relating `Why?' and `Why Not?' explanations.
\emph{arXiv preprint arXiv:2012.11067}, 2020.
\url{https://arxiv.org/abs/2012.11067}.

\bibitem[Kumar and Ahuja(2026)]{kumar2026}
Shubham Kumar and Narendra Ahuja.
Minimal, local, causal explanations for jailbreak success in large language
models.
In \emph{Conference on Language Modeling}, 2026.
\url{https://arxiv.org/abs/2605.00123v3}.

\bibitem[Li et~al.(2023)]{li2023}
Xuan Li, Zhanke Zhou, Jianing Zhu, Jiangchao Yao, Tongliang Liu, and Bo Han.
DeepInception: Hypnotize large language model to be jailbreaker.
\emph{arXiv preprint arXiv:2311.03191}, 2023.
\url{https://arxiv.org/abs/2311.03191}.

\bibitem[Zeller and Hildebrandt(2002)]{zeller2002}
Andreas Zeller and Ralf Hildebrandt.
Simplifying and isolating failure-inducing input.
\emph{IEEE Transactions on Software Engineering}, 28(2):183--200, 2002.
\url{https://doi.org/10.1109/32.988498}.

\bibitem[Zhao et~al.(2025)]{zhao2025}
Haiquan Zhao, Chenhan Yuan, Fei Huang, Xiaomeng Hu, Yichang Zhang,
An Yang, Bowen Yu, Dayiheng Liu, Jingren Zhou, Junyang Lin, and others.
Qwen3Guard technical report.
\emph{arXiv preprint arXiv:2510.14276}, 2025.
\url{https://arxiv.org/abs/2510.14276}.

\end{thebibliography}

\appendix
\section{Proofs and illustrative edge cases}
\label{app:proofs}

\begin{proof}[Proof of Proposition~\ref{prop:bounds}]
If $S\in\Nec(g)$, every completion has $f(S)=1$ and $f(T)=0$ for every
strict subset, so $S$ belongs to every minimal family. Conversely, if
$g(S)\ne1$, either it is already zero or a completion can set it to zero;
then $S$ is not necessary. If $g(S)=1$ but some $T\subsetneq S$ is not
known zero, a completion can set that $T$ to one, also preventing minimality.
This proves the intersection identity.

If $S$ is minimal in any completion, its label cannot be known zero and
no strict subset can be known one, which gives $S\in\Poss(g)$. Conversely,
for any $S\in\Poss(g)$, choose a completion that sets $f(S)=1$ and all
unresolved strict subsets to zero. This is consistent with every known label;
other unknown rows may be set arbitrarily. Then $S\in\Fam(f)$, proving
the union identity. A common intersection and union identify all memberships;
if they differ, the witnessing completions disagree on at least one member.
\end{proof}

\begin{proof}[Proof of Proposition~\ref{prop:closure}]
For monotone $f$, each $\cl_\pi(M)$ with $M\in\Fam(f)$ is positive.
Every positive representable set $A$ contains some fine minimum $M$, since
the finite collection of positive subsets of $A$ has a minimal member.
Representability of $A$ implies $\cl_\pi(M)\subseteq A$. Thus every
coarse positive contains a positive closure of a fine minimum. Its minimal
members are exactly the inclusion-minimal closures, proving the identity.
The argument also covers the empty family and a positive empty set.
\end{proof}

\paragraph{Why immediate subsets are insufficient.}
Let $f(\{a\})=1$, set every two-element subset of $\{a,b,c\}$ to zero,
and let $f(\{a,b,c\})=1$. All immediate predecessors of the three-element
set are negative, but the set is not strict-subset minimal because
$\{a\}$ is positive. This is a synthetic logical example, not an empirical
model response.

\paragraph{Why a zero residual does not prove monotonicity.}
For $U=\{a,b,c\}$, take positives $\{a\}$, $\{b\}$, and
$\{a,b,c\}$ only. The function is not monotone, since $\{a\}$ is
positive and $\{a,c\}$ is negative. For the partition with the single
block $U$, however, the coarse minimum and the minimal closure of both fine
minima are $U$. Agreement for a tested partition is therefore compatible
with nonmonotonicity elsewhere.

\section{Recorded observation contract}
\label{app:contract}

The D3 contract maps the Qwen3Guard axis to harmful when
\texttt{safety == Unsafe} and \texttt{refusal == No}; its non-success
predicate is \texttt{safety == Safe} or \texttt{refusal == Yes}. The panel
returns harmful if the Qwen harmful predicate and JailMeter label 1 agree,
safe if the Qwen non-success predicate and JailMeter label 0 agree, and
abstain otherwise, including parsing failure. Harmful takes precedence,
then safe, then abstain. The same recorded panel is used across the
scientific conditions. No evaluator output was changed in the retrospective
audit. This operating rule is not claimed to be a universal equivalence
between absence of attack success and safety.

Recovery requires safe panel labels for seeds 11, 23, and 47 under both
neutralizers, plus capability checks for the candidate recovery. Two benign
control tasks, arithmetic and token copying, are applied to separate control
copies of the prompt under each neutralizer. Their decoding uses seed 17,
temperature zero, top-$k$ one, top-$p$ one, and at most 48 new tokens.
Both tasks must return the required answer without truncation, prompt echo,
or runtime error. These four matched checks explain the 324 controls for
81 candidate-positive subsets. They provide a limited validity filter,
not a general capability benchmark.

Source-aware omission removes the designated source components through the
fixed renderer. Layout-preserving blanking replaces their owned regions
according to the frozen adapter. The current audit changes neither
neutralizer nor their definitions. The paired neutralizer-stability gate
compares their positive decisions before the joint recovery decision; its
pooled Jaccard therefore has a different denominator from the count of
jointly recovered subsets. The value $64/141$ is not a count of 141
independent payloads.

\section{Instance-level identification results}
\label{app:instances}

Table~\ref{tab:instances} preserves every instance, including the instance
without a necessary minimum. The safe hash prefixes identify the original
rows without reproducing request or response text. A possible count greater
than a necessary count marks an unidentified family.

\begin{table}[h]
\caption{All instance-level results. Unknown rows include both abstention
and capability confounding. ``Reversal contrasts'' counts adjacent partitions.}
\label{tab:instances}
\centering
\small
\begin{tabular}{llrrrr}
\toprule
ID & Safe instance prefix & Unknown & Necessary & Possible & Reversal contrasts \\
\midrule
D1 & \texttt{5233344c} & 25 & 2 & 10 & 2 \\
D2 & \texttt{8b1ff2f6} & 12 & 1 & 14 & 0 \\
D3 & \texttt{e40408e7} & 7 & 1 & 5 & 2 \\
H1 & \texttt{188c2aa2} & 0 & 1 & 1 & 0 \\
H2 & \texttt{27ac88c5} & 0 & 1 & 1 & 0 \\
H3 & \texttt{3b6bb0aa} & 0 & 1 & 1 & 0 \\
H4 & \texttt{7440b9c7} & 0 & 1 & 1 & 0 \\
H5 & \texttt{945bd2bd} & 1 & 2 & 3 & 0 \\
H6 & \texttt{bf3e47d0} & 0 & 1 & 1 & 0 \\
H7 & \texttt{cccc036f} & 1 & 0 & 3 & 0 \\
H8 & \texttt{d55a7656} & 0 & 3 & 3 & 0 \\
H9 & \texttt{de27373c} & 0 & 2 & 2 & 0 \\
\bottomrule
\end{tabular}
\end{table}

\section{Artifact identities and reconstruction}
\label{app:artifacts}

The new pure implementation is
\path{src/jbspan/topology_identification.py}. Its audit entry point is
\path{scripts/audit_rescue_identification_v2.py}; from the repository root,
\texttt{python scripts/audit\_rescue\_identification\_v2.py --root .}
reconstructs the results from the existing safe D3 source. It verifies the
source hash and reconstructs every historical certified minimum before
performing the new audit. The result is written once to
\path{data/natural_language_localization/rescue_identification_v2/audit.safe.json};
a later run with different content is rejected. No target, judge, payload,
or private-response access is needed by this entry point.

The D3 safe source is
\path{data/natural_language_localization/d3_exact_topology_v1/instance_results.safe.jsonl}.
Its SHA-256 is
\begin{center}\scriptsize
\texttt{2ebe3ff8cf7a603fa885500e9946bb79cdc7bd80d79f6ea29077dfb099934d34}.
\end{center}
The current audit result identity is
\begin{center}\scriptsize
\texttt{52e114d41d998e161c720b6ce89da87d8da908e89b2d7e3e873e3b0d8e102acc}.
\end{center}
The first output is retained separately as \path{audit.initial.safe.json};
its scientific values are unchanged in the current output, whose identity
also binds the subsequently formatted source bytes. Operator and seed
provenance are reconstructed in
\path{data/natural_language_localization/rescue_identification_v2/witness_provenance.safe.json},
which deduplicates exact atomic reversal pairs and binds their original safe
observations. This secondary audit does not infer unobserved seed outcomes.
The immutable D3 result file SHA-256 is
\begin{center}\scriptsize
\texttt{12e9289370c3db20423736103c5ca4f12666e4faf77c8a5b3eeeffdb0ce7e59b},
\end{center}
and the C1N result file SHA-256 is
\begin{center}\scriptsize
\texttt{acd2467e4c805fe6c3fe0b5f2ba7d416fe08a00760032166051430253c02b948}.
\end{center}
Their separate independent-verification receipts certify reconstruction of
the negative results; they do not certify a positive scientific hypothesis.
The current manuscript is a separate file from the earlier prospective
manuscript, which remains preserved. Bibliographic metadata and the limited
related-work statements were checked against primary publisher, proceedings,
or author-deposited records on 2026-09-05.

\end{document}
